\section{Vorwort}

{\ttfamily\color{BaseDark}

\noindent Liebe Leserinnen und Leser!

\bigskip

% Die  Arbeits-und
% Ausbildungsstandards für den Sanitätsdienst (AASS) sind
% mehrere Bände gegliedert: Bd. 1 enthält die Teile I, II und III und
% vermitteln allgemeines Basiswissen für den Sanitätsdienst.
% Band 2 beinhaltet erweitertes Wissen, welches für eine
% Betreuung von Patienten nach dem Stand der Wissenschaft
% hilfreich bzw. für die Betreuung von Notfallpatienten
% wichtig ist. Beide Bände beinhalten eine Band übergreifendem
% Appendix mit einem Indexverzeichnis und Literatur
% Verzeichnis. 
% 
% 
% Zusätzlich zu dem AASS-Hauptwerk existiert ein
% Kompendium –- vormals Maßnahmekatalog --, welcher den
% aktuellen Maßnahmekatalog sowie wichtige Informationen zumm
% nachschlagen, welche für den alltäglichen Dienstbetrieb
% relevant sind, beinhaltet. 
% 
% Für spezielle
% Anwendungsbereiche existieren so genannte Auszüge, welche
% schwerpunktmäßig bestimmte Themen behandeln, zum Beispiel
% Schwerpunkt Reanimation oder Schwerpunkt
% Patientenmanagement.





% \marginpar{\includegraphics[width=\linewidth]{./images/gabriel-sebastian-aass/avatar-defense.jpg}}
% \parpic[r][r]{\includegraphics[width=2cm]{./images/gabriel-sebastian-aass/avatar-defense.jpg}}
% TODO Vorwort fertig machen
Das vorliegende Werk ist das Ergebnis einer seit dem Jahr 2008 
andauernden Arbeit. 
%
In Zusammenarbeit mit dem
Ausbildungszentrum des ASB Florisdorf-Donaustadt
entstanden, 
ausgehend von den Erfordernissen
der entsprechenden Kurse, 
die ersten Versionen der {\AASS}.
%
Sie bildeten die Grundlage der Curricula 
und konnten sich schon früh, 
etappenweise und unter kontrollierten Bedingungen eingeführt, 
als Lern- und Ausbildungsunterlage bewähren.


Oft wurden wir dabei
gefragt \Speech{\enquote{Wann wird Euer Buch fertig sein?}} 
Nun, das wird es wohl nie sein. 
Die {\AASS} sind darauf ausgelegt,
ständig weiterentwickelt, gewartet und verbessert zu werden. 
%Das medizinische
%und organisatorische Umfeld wandelt sich ständig, sodass man es sich gar nicht
%leisten könnte, still zu stehen.
Mit jeder Ausgabe hoffen wir
auf wertvolle Anregungen und konstruktive Kritik durch die
Leser- und Anwenderschaft.


\bigskip

\noindent\textbf{Das Ziel des Projekts: \enquote*{Empowerment}}\par%
\noindent\mbox{}\marginpar{\includegraphics[width=\linewidth]{./images/teaser/affe-sama-ambas-schein_00800.jpg}

\FontSansNarrowFamily\small Algorithmenaffe oder
Fachpersonal?}%
% 
% \parpic[r][r]{\includegraphics[width=2cm]{./images/teaser/affe-sama-ambas-schein_00800.jpg}}
% 
% 
% 
% 
% 
%
Das zentrale Konzept der \T{Aufklärung} 
wird von Kant als die 
\Speech{\enquote{Befreiung von der selbst verschuldete
Unmündigkeit}} definiert. 
In diesem Geiste entstanden 
und entstehen die
{\AASS}: 
%
Sie sollen einerseits dem Sanitätspersonal, 
egal ob ärztlich oder
nicht-ärztlich, 
eine souveräne, 
überlegte 
und gleichzeitig
zielgerichtete Betreuung von Patienten 
und Notfallpatienten
ermöglichen. 
%
Andererseits ist es unser erklärtes Ziel, 
das Fachpersonal zum
eigenständigen Denken, 
Verstehen von Zusammenhängen 
und Hintergründen 
und kritischen Hinterfragen zu befähigen:
Gelebtes \enquote*{Empowerment}. 


\bigskip

\noindent\textbf{Geleitet, nicht gebunden}\par
\noindent Auch wenn es in
Mode gekommen ist, 
Abläufe zu standardisieren 
und
festzuschreiben, 
so ist dies keine taugliche
Herangehensweise, 
um komplexe Probleme in der Medizin 
oder
im Sanitätsdienst zu bewältigen. 
Zu komplex sind die
Voraussetzungen, 
und zu gering 
bzw. zu unangepasst an konkrete
Situationen ist das gesicherte Wissen, 
auf dem
festgeschriebene Prozesse basieren. 
%
Der menschliche Faktor
und individuelle Erfahrungen sind 
und bleiben der
entscheidende Faktor jeglichen Handelns. 
%
Kein Lehrbuch oder
standardisierter Prozess kann das ändern, 
sofern das Ziel
ein \enquote*{gutes} Ergebnis, 
und nicht das Erreichen einer
statistischen Kennziffer ist.
%
Es ist Ausdruck der Würde
und Wertschätzung gegenüber den im Sanitätsdienst tätigen 
Personen, 
ihnen einen umfassenden
Einblick in jene Welt zu gewähren, 
in der sie arbeiten. 


\bigskip

\noindent\textbf{Zeitloses Wissen}\par
\noindent Besonderen Wert haben wir dabei 
auf das Hintergrundwissen und relevante Randbereiche gelegt.
%
In einer Zeit,
in der das praxisbezogene Wissen oft innerhalb weniger Jahre bereits
wieder veraltet ist, 
bildet dieses Fundament die Chance für eine langjährige Tätigkeit 
und erlaubt selbstständiges,
regel\E{geleitetes} Arbeiten 
(anstatt strikten, regel\E{gebundenen} Handelns).


So sind der Titel dieses Werkes 
und darin enthaltene Handlungsempfehlungen 
nicht im Sinne eines standardisierten, 
rigiden Prozessmanagements zu verstehen,
sondern sollen einen \E{Basisstandard} bezüglich
medizinisch-sanitätsdienstlicher Versorgung, 
Wissen und
Kompetenz bilden, 
der individuell 
-- je nach Situation 
bzw. durch vom Helfer erworbenen Wissens -- 
jederzeit verbesser- und erweiterbar ist. 
%
Die {\AASS} geben Empfehlungen, 
es obliegt dem Helfer diese gemäß seiner 
fachlichen Erfahrung und Kenntnisse,
gepaart mit einer großen Prtion selbstkritischer Reflexion, 
zu interpretieren, 
einzusetzen 
oder anzupassen -- 


\Speech{salus aegroti suprema lex}\footnote{salus aegroti suprema lex \Lang{lat.}: \Speech{Das
Wohl des Kranken ist oberstes Gebot.}}.

\bigskip

\bigskip

\begin{flushright}

\emph{Wien, im Jänner 2015}

\bigskip

\emph{Die Herausgeber}
\end{flushright}

\normalfont
}